\section{Problem Statement}
This section introduces important terminologies and defines the problem of integrating anomaly detection and root cause localization with the same inputs.

\subsection{Terminologies}
Traces record the process of the microservice system responding to a user request (e.g., click ``create an order'' on an online shopping website). 
Different microservice instances then conduct a series of small actions to respond to the request.
For example, the request ``create an order'' may contain steps ``create an order in pending'', ``reserve credit'', and ``update the order state.''
A microservice (caller) can \textit{invoke} another microservice (callee) to conduct the following action (e.g., microservice ``Query'' asks microservice ``Check'' to check the order after finishing the action ``query the stock of goods''), and the callee will return the result of the action to the caller. We name this process as \textit{invocation}.
The time consumed by the whole invocation (i.e., from initializing the invocation to returning the result) is called invocation \textit{latency}, including the request processing time inside a microservice and the time spent on communicating between the caller and the callee. 
A \textit{trace} records the information during processing a user request~\cite{MEPFL}  (including multiple invocations), such as the invocation latency, the total time of processing the request, the HTTP response code, etc.

Meanwhile, system logs are generated when system events are triggered. 
A \textit{log message} (or \textit{log} for short) is a line of the standard output of logging statements, composed of constant strings (written by developers) and variable values (determined by the system)~\cite{SurveyHe}. If the variable values are removed, the remaining constant strings constitute a \textit{log event}. 
\textit{KPIs} are the numerical measurements of system performance (e.g., disk I/O rate) and the usage of resources (e.g., CPU, memory, disk) that are sampled uniformly. 

\subsection{Problem Formulation}\label{sec:prob form}
Consider a large-scale system with $M$ microservices, system logs, KPIs, and traces are aggregated individually at each microservice.
In a $T$-length observation window (data obtained in a window constitute a sample), we have multi-source data defined as $\mathbf{X}=\{ ( \mathbf{X}^\mathcal{L}_{m}, \mathbf{X}^\mathcal{K}_{m}, \mathbf{X}^\mathcal{T}_{m} ) \}_{m=1}^M$, 
where at the $m$-th microservice,
$\mathbf{X}^\mathcal{L}_{m}$ represents the log events chronologically arranged;
$\mathbf{X}^\mathcal{K}_{m}$ is a multivariate time series consisting of $k$ indicators;
$\mathbf{X}^\mathcal{T}_{m}$ denotes the trace records.
Our work attempts to build an end-to-end framework achieving a two-stage goal:
Given $\mathbf{X}_{[1:M]}$, the framework predicts the existence of anomalies, denoted by $y$, a binary indicator represented as 0 (normal) or 1 (abnormal). 
If $y$ equals one, a localizer is triggered to estimate the probability of each microservice to be the culprit, denoted by $\textbf{P}=[p_{1} \cdots p_{M}] \in [0, 1]^M$.
The framework is built on a parameterized model $\mathcal{F}: \mathbf{X} \rightarrow (y, \mathbf{P})$.
